\documentclass[10pt]{article}
% Math Packages
\usepackage{amsmath, mathtools}
\usepackage{amssymb}
\usepackage{amsthm}
\usepackage{amsfonts}
\usepackage{bbm}
\usepackage{breqn}
\usepackage[margin=1in]{geometry}
\usepackage{graphicx}
\usepackage{tikz}
\usetikzlibrary{arrows.meta}
\usetikzlibrary{calc}
\usepackage{forest}
\usepackage{tikz-qtree}
\graphicspath{ {./images/} }
\usepackage{hyperref}
\usepackage[capitalize]{cleveref}
\usepackage[shortlabels]{enumitem}
\usetikzlibrary{arrows,matrix,positioning}
\usepackage{multicol}

% for the pipe symbol
\usepackage[T1]{fontenc}

% Citing theorems by name. (source: https://tex.stackexchange.com/questions/109843/cleveref-and-named-theorems)
\makeatletter
\newcommand{\ncref}[1]{\cref{#1}\mynameref{#1}{\csname r@#1\endcsname}}

\def\mynameref#1#2{%
  \begingroup
  \edef\@mytxt{#2}%
  \edef\@mytst{\expandafter\@thirdoffive\@mytxt}%
  \ifx\@mytst\empty\else
  \space(\nameref{#1})\fi
  \endgroup
}
\makeatother

% Colorful Notes
\usepackage{color} \definecolor{Red}{rgb}{1,0,0} \definecolor{Blue}{rgb}{0,0,1}
\definecolor{Purple}{rgb}{.5,0,.5} \def\red{\color{Red}} \def\blue{\color{Blue}}
\def\gray{\color{gray}} \def\purple{\color{Purple}}
\newcommand{\rnote}[1]{{\red [#1]}} % \rnote{foo} gives '[foo]' in red
\newcommand{\pnote}[1]{{\purple [#1]}} % \pnote{foo} gives '[foo]' in purple
\newcommand{\bnote}[1]{{\blue #1}} % \bnote{foo} gives 'foo' in blue
\newcommand{\gnote}[1]{{\gray #1}} % \gnote{foo} gives 'foo' in gray
\newcommand{\Max}[1]{{\purple [#1]}} % \bnote{foo} then 'foo' is blue


% Claim numbering (the counter restarts after each proof environment)
\newcounter{claimcount}
\setcounter{claimcount}{0}
\newenvironment{claim}{\refstepcounter{claimcount}\par\addvspace{\medskipamount}\noindent\textbf{Claim \arabic{claimcount}:}}{}
\usepackage{etoolbox}
\AtBeginEnvironment{proof}{\setcounter{claimcount}{0}}
\newenvironment{claimproof}{\par\addvspace{\medskipamount}\noindent\textit{Proof of Claim  \arabic{claimcount}.}}{\hfill\ensuremath{\qedsymbol} \tiny{Claim}

  \medskip}
% Add claim support to cleverref
\crefname{claimcount}{Claim}{Claims}

% Math Environments
\newtheorem{theorem}{Theorem}
\newtheorem{assumption}[theorem]{Assumption}
\newtheorem{lemma}[theorem]{Lemma}
\newtheorem{proposition}[theorem]{Proposition}
\newtheorem{corollary}[theorem]{Corollary}
\newtheorem{question}[theorem]{Question}
\theoremstyle{definition}
\newtheorem{definition}[theorem]{Definition}
\newtheorem{remark}[theorem]{Remark}
\newtheorem{example}[theorem]{Example}
\newtheorem{notation}[theorem]{Notation}
\newtheorem{problem}[theorem]{Problem}

% Matrices and Column Vectors. 
\usepackage{stackengine}
\setstackgap{L}{1.0\normalbaselineskip}
\usepackage{tabstackengine}
\setstackEOL{;}% row separator
\setstackTAB{,}% column separator
\setstacktabbedgap{1ex}% inter-column gap
\setstackgap{L}{1.5\normalbaselineskip}% inter-row baselineskip
\let\nmatrix\bracketMatrixstack  %Usage: \nmatrix{1,2,3\4,5,6}
\newcommand\cv[1]{\setstackEOL{,}\bracketMatrixstack{#1}} %usage: \cv{1,2,3}

% Custom Math Coqmmands
\newcommand{\vt}{\vskip 5mm} % vertical space
\newcommand{\fl}{\noindent\textbf} % first line
\newcommand{\Fl}{\vt\noindent\textbf} % first line with space above
\newcommand{\norm}[1]{\left\lVert#1\right\rVert} % norm
\newcommand{\pnorm}[1]{\left\lVert#1\right\rVert_p} % p-norm
\newcommand{\qnorm}[1]{\left\lVert#1\right\rVert_q} % q-norm
\newcommand{\1}[1]{\textbf{1}_{\left[#1\right]}} % indicator function 
\def\limn{\lim_{n\to\infty}} % shortcut for lim as n-> infinity
\def\sumn{\sum_{n=1}^{\infty}} % shortcut for sum from n=1 to infinity
\def\sumkn{\sum_{k=1}^{n}} % shortcut for sum from k=1 to n
\def\sumin{\sum_{i=1}^{n}} % shortcut for sum from i=1 to n
\def\SAs{\sigma\text{-algebras}} % shortcut for $\sigma$-algebras
\def\SA{\sigma\text{-algebra}} % shortcut for $\sigma$-algebra
\def\Ft{\mathcal{F}_t} % time-indexed sigma-algebra (t)
\def\Fs{\mathcal{F}_s} % time-indexed sigma-algebra (s)
\def\F{\mathcal{F}} % sigma-algebra
\def\G{\mathcal{G}} % sigma-algebra
\def\R{\mathbb{R}} % Real numbers
\def\N{\mathbb{N}} % Natural numbers
\def\Z{\mathbb{Z}} % Integers
\def\E{\mathbb{E}} % Expectation
\def\P{\mathbb{P}} % Probability
\def\Q{\mathbb{Q}} % Q probability
\def\dist{\text{dist}} %Text 'dist' for things like 'dist(x,y)'
\newcommand{\indep}{\perp \!\!\! \perp}  %independence symbol
\def\Var{\mathrm{Var}} % Variance
\def\tr{\mathrm{tr}} % trace

% Brackets and Parentheses
% \def\[{\left [}
%     \def\]{\right ]}
% \def\({\left (}
%   \def\){\right )}



\usepackage{color}
\definecolor{Red}{rgb}{1,0,0}
\definecolor{Blue}{rgb}{0,0,1}
\definecolor{Purple}{rgb}{.5,0,.5}
\def\red{\color{Red}}
\def\blue{\color{Blue}}
\def\gray{\color{gray}}
\def\purple{\color{Purple}}
\definecolor{RoyalBlue}{cmyk}{1, 0.50, 0, 0}

\DeclareRobustCommand{\demph}[1]{% definitions
  \ifmmode
  {\color{Blue}#1}%
  \else
  {\color{Blue}\bfseries\slshape #1}%
  \fi
}
% comment exactly one of the following line to show / hide the solutions
% \newcommand{\solution}[1]{{\purple #1}} % uncomment to show the solutions
\newcommand{\solution}[1]{} % uncomment to hide the solutions



\title{Lecture Notes for Math 372: \\Elementary Probability and Statistics}
\date{Last updated: \today}
% \author{mh}

\begin{document}
\begin{center}
  \section*{Math 243: Homework 03}
  \textit{Due Fridy, September 11}
\end{center}


\begin{problem}
  Consider
  \[
    \mathbf r(t)=(\cos(2t),\sin(2t)),
    \qquad 0\leq t\leq2\pi.
  \]
  This parameterization traces out the unit circle, whose circumference is
  $2\pi$, but $\int_0^{2\pi}|\mathbf r'(t)|\,dt=4\pi.$
  Explain why there is no contradiction. What does the integral measure in
  this example?
\end{problem}

\begin{problem}
  Consider $\mathbf r(t)=(t^{3},t^{2})$.
  \begin{enumerate}[(a)]
    \item Compute the velocity at time $t=0$. Why does the direction formula
    $\mathbf r'(t)/\left|\mathbf r'(t)\right|$ fail at $t=0$?

    \item Convert the curve to Cartesian form by eliminating the variable $t$
    and sketch the curve. Explain geometrically what goes wrong.
  \end{enumerate}
\end{problem}

\begin{problem}[Converting parametric equations to Cartesian form]
  Consider the curve given by the parametric equations
  \begin{equation*}
    x(t) = \frac{t-2}{t+1}, \quad y(t) = \frac{t}{t-1}
  \end{equation*}
  where  $-\frac{1}{2}\leq t \leq \frac{1}{2}$. 

  \begin{enumerate}[(a)]
    \item Find a Cartesian equation for the curve. Sketch a graph, and
    show the direction of travel. \textit{(Hint: be careful about precisely which values of $x$ your equation
    holds for.)}
    \item How does your answer change if we specify $-1<t<1$? Plot the graph in that case as well. \textit{(Note
    that we can't take $t=-1$ or $t=1$ because that would result in divison by zero!)}
\end{enumerate}

\end{problem}


\begin{problem}[Arc length of parameterized curve]
  Compute the arc length of the parameterized curve:
  \begin{equation*}
    \left\{ \begin{array}{l@{}l} x(t) = \frac{4}{3}t^{3}+t-2 \\ y(t) = \frac{1}{27}\left( 9t^{2}+2 \right)^{\frac{3}{2}} 
      \end{array}\right.
  \end{equation*}
  for $0 \leq t \leq 1$. (Don't forget to use the chain rule!)
\end{problem}


\begin{problem}[Arc length - 3d]
  Find the length of the following curve in 3d space:
  \begin{equation*}
    \mathbf{r}(t)=
    \begin{bmatrix}
      x(t)\\y(t)\\z(t)
    \end{bmatrix}
    =
    \begin{bmatrix}
      1\\t^{2}\\t^{3}
    \end{bmatrix}
    , \quad 0 \leq t \leq 9.
  \end{equation*}
\end{problem}

\begin{problem}[Parametric equations - tangent lines]
  Let $C$ be the curve defined by the parametric equations
  \begin{equation}
    \left\{ \begin{array}{l@{}l} x(t) = t^{3}-12t \\
        y(t) = 10-t^{2} \end{array}\right.
  \end{equation}
  for $0 \leq t \leq 3$.

  
  Here's a plot of the graph of this curve: \url{https://www.desmos.com/calculator/bvti5ck5r0}
  
  \begin{enumerate}[(a.)]
    \item Find an equation for the tangent line to $C$ at $t=1$. Your final answer
    should be an equation of a line in \textit{slope-intercept form:} $y=mx+b$.
    \item What is the equation of the normal line to $C$ at $t=1$?
    \item Sketch the curve, the tangent line, and the normal line. (Please
    label them).
    \item Compute the value of $\frac{d^{2}y}{dx^{2}}$ at the point $t=1$.
    \item Based on your answer in part (d.), is the curve $C$ concave up or
    concave down at $t=1$? \textit{(Note: You can get the answer by looking at
      the graph, but the point of this part is that you know how to interpret
      part (d.)}
    \vspace{3em}
  \end{enumerate}
 
  
  \begin{problem}[Differentiating parametric equations]
    Consider the curve $C$ given by the parametric equations
    \begin{equation*}
      \left\{ \begin{array}{l@{}l} x=\frac{t}{1+t} \\ y = \sqrt{1+t} \end{array}\right.
    \end{equation*}
    for $0 \leq t <\infty$.
    \begin{enumerate}[(a.)]
      \item Compute $\frac{dy}{dx}.$
      \item Let $f(t)=\frac{dy}{dx}(t)$. In your words, provide a geometric interpretation or explanation
      of the function $f$.
      \item Compute the second derivative $\frac{d^{2}y}{dx^{2}}$.
      \item Based on your answer in part (c.), at the point $t=1$, is the curve $C$ concave up, concave down, or
      neither? Justify your answer.
    \end{enumerate}
  \end{problem}

\end{problem}



% \newpage
% \subsection{Unused}



\end{document}